% Generated draft: requires mathematical and source review.
\documentclass[11pt]{article}
\usepackage[margin=25mm]{geometry}
\usepackage{fontspec}
\setmainfont{Latin Modern Roman}
\usepackage{amsmath,amssymb,mathtools,unicode-math}
\setmathfont{Latin Modern Math}
\usepackage{enumitem,needspace,xurl,hyperref}
\setlength{\parindent}{0pt}
\setlength{\parskip}{5pt}
\setlength{\emergencystretch}{3em}
\newcommand{\N}{\mathbb{N}}
\newcommand{\Z}{\mathbb{Z}}
\newcommand{\Q}{\mathbb{Q}}
\newcommand{\R}{\mathbb{R}}
\newcommand{\C}{\mathbb{C}}
\newcommand{\F}{\mathbb{F}}
\newcommand{\A}{\mathbb{A}}
\newcommand{\PP}{\mathbb{P}}
\newcommand{\E}{\mathbb{E}}
\newcommand{\eps}{\varepsilon}
\newcommand{\dd}{\,\mathrm{d}}
\begin{document}
\section*{UnsolvedMath ID: 11000118}
\textbf{Source problem code:} catalogue entry
\subsection{Definitions and mathematical statement}
All homology groups in this problem have integral coefficients. Retain the source's \(\operatorname{Mod}_g\), \(\operatorname{Mod}_g^1\) and pure-group notation \(\operatorname{PMod}_{g,p}^{b}\); the meanings of the superscript and the two marking/boundary indices are not defined in the extract.
\begin{enumerate}
\item In genus \(3\), the source reports from [33] that each of
\[
 H_2(\operatorname{Mod}_3;\Z),\qquad
 H_2(\operatorname{Mod}_3^1;\Z)
\]
is isomorphic either to \(\Z\) or to \(\Z\oplus\Z/2\Z\). Determine which group occurs in each case. Also compute
\[
 H_2(\operatorname{PMod}_{3,p}^{b};\Z)
 \qquad\text{for every allowed }p,b.
\]
\item The source reports
\(H_2(\operatorname{Mod}_2;\Z)\cong H_2(\operatorname{Mod}_2^1;\Z)\cong\Z/2\Z\).
Compute \(H_2(\operatorname{PMod}_{2,p}^{b};\Z)\) for \(b\geq2\) and all allowed \(p\).
\end{enumerate}
For a pure mapping class group the usual requirement is fixing punctures individually, and any boundary convention must also be specified. Assigning \(b\) versus \(p\) to punctures or boundary components without the omitted definitions could change the problem. The computation asks for the full second homology group, including torsion, not only its rational rank.

The supplied research summary says the initial genus-three torsion alternatives have been resolved in later literature, but explicitly lacks an independently verified reference and gives no exact answer. That partial-resolution report is retained without inventing a value or treating all of the two-parameter computations as solved.
\subsection{Short English statement}
Resolve the stated genus-three second-homology torsion alternatives and compute the full integral second homology of the indicated pure groups in genera two and three, retaining the exact parameter ranges.
\subsection{Sources}
\begin{itemize}
\item \emph{Problems on Mapping Class Groups and Related Topics}, edited by Farb, Problem 2.13, supplied PDF page 98. \url{https://www.math.uchicago.edu/~farb/papers/mcgbook.pdf}. Recorded access 2026-07-29; extraction pdf-font-structure.
\item Source reference [33] for the genus-three alternatives; its bibliography entry is absent. The summary loosely mentions Harer--Ivanov/cohomology computations but supplies no exact later theorem or reference for the claimed resolution.
\item Local UnsolvedMath registry, \texttt{data/unsolvedmath/problems.json}, permanent ID \(11000118\), exact source code \texttt{catalogue entry}; registry status open/OPEN-TRIAGE, with a partial-resolution report in the summary.
\end{itemize}
\end{document}
